% !TEX root = ../main.tex


\section{The Clearing Rule}
\label{app:zq-clearing-rule}

We define \(\mathsf{Clear}\) on books with ordinary and imbalance-only orders. Its output is either a clearing receipt or a no-cross outcome.

\subsection{Definition}

A \emph{book} on \(n\) ticks, of which the first \(N\) are real, is four sequences of integers \(b_i, a_i, b^{io}_i, a^{io}_i\), \(0 \le i < n\), with every entry in \([0, 2^{k_{in}})\) and zero for \(i \ge N\). \(b_i\) is the bid volume with limit exactly tick \(i\) and \(a_i\) the ask volume; a market buy is recorded at tick \(N-1\) and a market sell at tick 0. The \(io\) sequences hold the imbalance-only orders. For each tick define

\[
D_i = \sum_{j \ge i} b_j, \quad
S_i = \sum_{j \le i} a_j, \quad
V_i = \min(D_i, S_i), \quad
g_i = S_i - D_i, \quad
\delta_i = |g_i|,
\]

and likewise \(D^{io}_i = \sum_{j \ge i} b^{io}_j\) and \(S^{io}_i = \sum_{j \le i} a^{io}_j\). \(D_i\) is the demand willing to trade at price \(i\), \(S_i\) the supply, and \(V_i\) the volume that can execute. \(D\) is non-increasing; \(S\) and \(g\) are non-decreasing. Tick \(i\) is \emph{demand long} if \(g_i < 0\), \emph{supply long} if \(g_i > 0\), and \emph{balanced} if \(g_i = 0\). Its \emph{at-price long volume} is \(\lambda_i = b_i\) if it is demand long and \(\lambda_i = a_i\) if it is supply long.

\textbf{Definition 2 (\(\mathsf{Clear}\)).} Let \(V = \max_i V_i\). If \(V = 0\), \(\mathsf{Clear}(\text{book}) = \textsf{no-cross}\). Otherwise let

\begin{itemize}
\item
  \(\Pi = \{i : V_i = V\}\) be the \emph{plateau}, with \(c = \min \Pi\) and \(d = \max \Pi\);
\item
  \(v = \min_{i \in \Pi} \delta_i\) be the minimum imbalance;
\item
  a \emph{candidate} be a tick \(i \in \Pi\) with \(\delta_i = v\) and (\(v = 0\) or \(\lambda_i \ge v\));
\item
  \(p^*\) be the least candidate.
\end{itemize}

Then \(\mathsf{Clear}(\text{book}) = (V, c, d, p^*, v, s, P, J)\), where \(s = P = J = 0\) if \(v = 0\), and otherwise \(s = [g_{p^*} > 0]\), \(P = \lambda_{p^*}\), and \(J = S^{io}_{p^*}\) if \(s = 0\), \(J = D^{io}_{p^*}\) if \(s = 1\).

\subsection{Properties}

\textbf{The plateau is an interval.} For \(i < j < l\) we have \(V_j \ge \min(V_i, V_l)\), because \(D_j \ge D_l \ge V_l\) and \(S_j \ge S_i \ge V_i\). Hence \(\Pi = \{c, \ldots, d\}\).

\textbf{Existence.} Lemma~\ref{app:zq-security}.1 supplies a candidate for every crossing book. Together with the no-cross case, this makes \(\mathsf{Clear}\) a total function.

\textbf{At most two candidates when \(v > 0\).} On \(\Pi\), \(g\) is non-decreasing and takes the values \(-v\) and \(+v\) on the minimizers. Let \(j\) be a demand-long minimizer that is not the highest one. Then \(g_{j+1} = g_j\), and since \(g_{j+1} - g_j = a_{j+1} + b_j\), we get \(b_j = 0 < v\): \(j\) is not a candidate. By symmetry only the lowest supply-long minimizer can be a candidate. The two candidates are therefore the highest demand-long minimizer and the lowest supply-long minimizer, and the first, when it exists, is the lower tick. When \(v = 0\) every balanced plateau tick is a candidate and \(p^*\) is the lowest.

\textbf{At-price volume.} We require better-priced orders to fill in full even when no IO volume is available. In that case all \(v\) unfilled units must come from long-side interest at the clearing tick, so we require \(\lambda_{p^*}\ge v\) when choosing the price. Proposition~\ref{app:zq-security}.4 shows that this condition yields feasible fill ratios for every eligible IO volume.

\subsection{Fills}
\label{app:zq-fills}

Of the eligible IO volume \(J\), the exchange admits \(I=\min(J,v)\), leaving residual imbalance \(R=v-I\). Both sides trade \(V+I\) units before rounding. Table 8 gives the pro rata entitlement for an order of size \(q\). When \(v=0\), all eligible ordinary orders fill in full and IO orders receive zero. A ratio is used only when its denominator is positive. The venue specifies integer rounding separately.

\textbf{Table 8.} The fill of an order of quantity \(q\), given the receipt.

\begin{longtable}{@{}
  >{\raggedright\arraybackslash}p{\dimexpr0.75\linewidth-\tabcolsep\relax}
  >{\raggedright\arraybackslash}p{\dimexpr0.25\linewidth-\tabcolsep\relax}
@{}}
\toprule
Order & Fill \\
\midrule
\endhead
\bottomrule
\endlastfoot

Limit or market order strictly better than \(p^*\), or on the short side at \(p^*\)
  & \(q\) \\
Limit order at exactly \(p^*\) on the long side
  & \(q \cdot (P - R)/P\) \\
Imbalance-only order on the short side, eligible at \(p^*\)
  & \(q \cdot I/J\) \\
Any other order
  & 0 \\
\end{longtable}

\subsection{The Running Example}

\textbf{Table 9.} The book of Figure 2, tick by tick.

\begin{longtable}[]{@{}lllllllll@{}}
\toprule\noalign{}
Tick \(i\) & 0 & 1 & 2 & 3 & 4 & 5 & 6 & 7 \\
\midrule\noalign{}
\endhead
\bottomrule\noalign{}
\endlastfoot
\(b_i\) & 0 & 0 & 3 & 0 & 0 & 10 & 2 & 0 \\
\(a_i\) & 0 & 4 & 6 & 0 & 0 & 0 & 5 & 3 \\
\(a^{io}_i\) & 0 & 0 & 0 & 0 & 1 & 0 & 0 & 0 \\
\(D_i\) & 15 & 15 & 15 & 12 & 12 & 12 & 2 & 0 \\
\(S_i\) & 0 & 4 & 10 & 10 & 10 & 10 & 15 & 18 \\
\(V_i\) & 0 & 4 & 10 & 10 & 10 & 10 & 2 & 0 \\
\(g_i\) & \(-15\) & \(-11\) & \(-5\) & \(-2\) & \(-2\) & \(-2\) & 13 & 18 \\
\end{longtable}

At ticks 2 through 5 the executable volume is 10, and the respective imbalances are 5, 2, 2, and 2. Demand is long throughout. Of the three minimum-imbalance ticks, only tick 5 has bids resting at the tick itself. Its 10 units cover the two-unit imbalance, so \(p^*=5\), \(s=0\), and \(P=10\). The eligible sell IO contributes one unit (\(J=I=1\)), leaving \(R=1\). Ordinary asks sell 10 units and the IO sells one. Bids at tick 6 buy two units and bids at tick 5 buy nine, for 11 on each side.

\subsection{Relation to the Nasdaq Closing Cross}

Nasdaq Rule 4754(b)(2) first maximizes executed shares, then minimizes imbalance. For remaining ties it considers the entered price at which shares stay unexecuted and finally proximity to the continuous market's bid--ask midpoint [37, Rule 4754(b)(2)]. Table 10 gives the corresponding steps in our rule.

\textbf{Table 10.} The steps of the Closing Cross and of \(\mathsf{Clear}\).

\begin{longtable}{@{}
  >{\raggedright\arraybackslash}p{\dimexpr0.32\linewidth-\tabcolsep\relax}
  >{\raggedright\arraybackslash}p{\dimexpr0.36\linewidth-2\tabcolsep\relax}
  >{\raggedright\arraybackslash}p{\dimexpr0.32\linewidth-\tabcolsep\relax}
@{}}
\toprule
Closing Cross & \(\mathsf{Clear}\) & Difference \\
\midrule
\endhead
\bottomrule
\endlastfoot

(A) maximize executed shares
  & \(V\) and the plateau \(\Pi\)
  & None \\

(B) minimize the imbalance
  & \(v\) and the minimizers
  & None \\

(C) the entered price at which shares remain unexecuted
  & Candidate: the long side's volume at the tick covers the imbalance
  & We state the step as the inequality \(\lambda_i \ge v\) \\

(D) closest to the bid-ask midpoint
  & The least candidate
  & A batch market has no continuous quote to refer to \\
\end{longtable}

The Closing Cross can draw on the continuous book. Our batch contains only orders represented in its four columns. We admit IO orders after setting the price, so they affect execution volume alone. At the cross price, Nasdaq Rule 4754(b)(3) allocates by order type and time priority. We use pro rata allocation on the long side, consistent with the absence of time priority within an FBA batch \cite{budish2015}.


